\chapter{Conclusion}
\label{cha:conclusion}

This study addressed two gaps in dictionary learning for CT-based CAD: the lack of a framework unifying
multi-class classification and voxel-level localisation, and the lack of any study benchmarking dictionary
learning against deep learning across classification, localisation, and explainability simultaneously. Four
dictionary learning methods (K-SVD, LC-KSVD2, FDDL, Frozen Dictionary Learning) were implemented in an
open-source package and evaluated, alongside three ML baselines and four SotA deep learning models, using a
shared, model-agnostic evaluation framework.

LC-KSVD outperformed all other dictionary learning models, reaching 88.5\% patch accuracy. Patch-level accuracy did not
transfer to whole-volume localisation for any dictionary learning method, a gap
driven by the extreme size imbalance between lesions and full volumes rather than by the sparse codes
themselves. Dictionary learning trails SotA deep learning substantially on raw performance, confirming a
interpretability-performance trade-off.

The study only covers two of five radiologist-recommended abnormalities and excludes the ambiguous
atelectasis/consolidation category. The fixed $12\times12\times6$ patch size discarded most candidate nodule
patches and biased the retained set toward larger masses.

Priority directions: spatially-aware post-processing to close the patch-to-volume localisation gap; extending
to the remaining radiologist-recommended abnormalities, including a properly disambiguated
atelectasis/consolidation split; domain-specific hyperparameter tuning (particularly LC-KSVD2's $\alpha/\beta$)
against a held-out validation set; a second independent radiologist review with inter-rater reliability
reporting, ideally with refined ground-truth masks; and confidence intervals or paired significance tests over
the closer method comparisons in Chapter~4.

Dictionary learning is not yet competitive with state-of-the-art deep learning on raw classification or
localisation accuracy for pulmonary abnormality detection, but the interpretability it offers is qualitatively
different from, and more reliable than, Grad-CAM's post-hoc
explanations on a leading vision-language model. A CAD system built on class-attributable atoms can show
exactly which learned pattern drove a prediction and where it is located.
The open-source benchmarking framework developed here is intended to make
that trade-off measurable for future work.
